\section{A3C: ein zweiter Algorithmus, eine Zerlegung und ein generatives Marktmodell}
\label{sec:a3c}

\paragraph{Woher die Frage kommt.} Original-Notebook und Nachbau (Abschnitt~\ref{sec:a2c})
nutzen \emph{A2C}, also synchrones Lernen. Die naheliegende Gegenprobe ist \emph{A3C}:
mehrere Worker sammeln parallel Rollouts und schieben ihre Gradienten in ein gemeinsames
Netz (Hogwild-Stil, nur der Optimizer-Schritt ist gesperrt). Umgebung, Fenster,
Episoden-Seeds und Auswertungsfunktionen bleiben unverändert --- damit misst der Vergleich
den Lernalgorithmus und nicht das Drumherum. Die Endauswertung läuft wieder über
\textbf{20 Seeds (100--119)}. Eine Besonderheit vorweg: A3C wurde erst
\emph{nach} der Fertigstellung des A2C-Teils gebaut, es gibt also keine
Notebook-Vorlage, gegen die man die Zahlen spiegeln könnte.

\subsection{Die Baseline hat keinen Vorsprung}

\begin{table}[H]
  \centering
  \caption{A3C auf Boeing, 20 Seeds. Marge $=$ Agent $-$ Buy-and-Hold in NAV-Punkten
           (Start $=1{,}0$); $p$ einseitig aus dem $t$-Test ``Marge $>0$''.
           ``R/R'' ist das Rendite/Risiko-Verhältnis im Testfenster.}
  \label{tab:a3c-arms}
  \small
  \begin{tabular}{lrrrrrr}
    \toprule
    Arm & Training & Test & 95-\%-KI Test & Treffer Test & $p$ & R/R Test\\
    \midrule
    Baseline                 & $+8{,}43$ & $+0{,}019$ & $[-0{,}06;+0{,}10]$ & 51,0\,\% & $0{,}32$  & $-0{,}24$\\
    \texttt{ent\_coef} 0,02  & $+8{,}37$ & $+0{,}047$ & $[-0{,}03;+0{,}12]$ & 64,0\,\% & $0{,}10$  & $-0{,}37$\\
    \texttt{obs\_noise} 0,01 & $+8{,}79$ & $+0{,}038$ & $[-0{,}03;+0{,}11]$ & 57,2\,\% & $0{,}14$  & $-0{,}45$\\
    \texttt{nstep} 10        & $+6{,}71$ & $+0{,}125$ & $[+0{,}01;+0{,}24]$ & 67,8\,\% & $0{,}019$ & $-0{,}05$\\
    Bündel (alle drei)       & $+7{,}69$ & $+0{,}145$ & $[+0{,}03;+0{,}26]$ & 70,2\,\% & $0{,}007$ & $-0{,}03$\\
    Bündel $+$ Bootstrap     & $+0{,}02$ & $+0{,}218$ & $[+0{,}09;+0{,}35]$ & 77,8\,\% & $0{,}0012$& $+0{,}50$\\
    \texttt{nstep} 10 $+$ Bootstrap & $+0{,}02$ & $+0{,}218$ & $[+0{,}09;+0{,}35]$ & 72,5\,\% & $0{,}0010$& $+0{,}38$\\
    \bottomrule
  \end{tabular}
\end{table}

Die Baseline wiederholt das A2C-Ergebnis: Im Testfenster liegt die Marge bei
$+0{,}019$ NAV-Punkten, der Agent gewinnt in \textbf{51,0\,\%} der Episoden --- ein
Münzwurf --- und der einseitige $p$-Wert ist $0{,}32$. \textbf{Ein A3C-Agent ohne
Anpassungen hat keinen belegbaren Vorsprung vor Buy-and-Hold.} Im Trainingsfenster steht
gleichzeitig eine Marge von $+8{,}43$: Der Agent hat den einen historischen Pfad
auswendig gelernt.

\subsection{Zerlegung: welcher Faktor trägt?}

Ein naheliegender Griff ins Werkzeugregal ist ein Bündel aus drei Änderungen:
\texttt{nstep} $5\to10$ (n-Schritt-Returns), \texttt{ent\_coef} $0{,}01\to0{,}02$
(Entropie-Bonus) und \texttt{obs\_noise} $0\to0{,}01$ (Gauß-Rauschen auf die Beobachtung
während des Trainings). Das Bündel hebt die Testmarge auf $+0{,}145$
($p=0{,}007$). Nur: \emph{welcher} der drei Faktoren tut das?

Deshalb wurde jeder Faktor einzeln bei ebenfalls 20 Seeds gefahren (Zeilen 2--4 in
Tabelle~\ref{tab:a3c-arms}). Das Ergebnis ist eindeutig und war für mich überraschend:

\paragraph{\texttt{nstep} ist der Wirkstoff.} Allein erreicht \texttt{nstep} $10$
bereits $+0{,}125$ ($p=0{,}019$, Trefferquote 67,8\,\%) --- das sind \textbf{85\,\%} der
Bündelwirkung. Der gepaarte Test Bündel gegen \texttt{nstep} allein ergibt
$+0{,}020$ bei $p=0{,}81$: statistisch \textbf{nicht unterscheidbar}. Das Bündel ist
also keine Summe von Hebeln, sondern im Wesentlichen \texttt{nstep} plus Rauschen.

\paragraph{\texttt{ent\_coef} und \texttt{obs\_noise} sind inert.} Einzeln liefern sie
$+0{,}047$ ($p=0{,}64$) und $+0{,}038$ ($p=0{,}72$) gegen die Baseline, mit
10 bzw.\ 9 von 20 positiven Seed-Differenzen --- ein Münzwurf. Auch auf das Overfitting
wirken sie nicht: Die In-Sample-Marge bleibt bei $+8{,}37$ bzw.\ $+8{,}79$ gegenüber
$+8{,}43$ in der Baseline ($p=0{,}90$ und $0{,}53$). Nur \texttt{nstep} $10$ senkt sie
signifikant auf $+6{,}71$ ($p=0{,}0015$) --- und hebt gleichzeitig den Testwert.

\paragraph{Additivität.} Die Summe der drei Einzel-Deltas ist $+0{,}155$, das Bündel
liegt bei $+0{,}126$; das Residuum ist $-0{,}029$ bei einer Streuung von $0{,}61$
($p=0{,}84$). Der Punktschätzer ist also leicht sub-additiv ($0{,}82\times$), aber von
reiner Additivität statistisch nicht zu unterscheiden --- das Rauschen dominiert.
Eine Synergie gibt es jedenfalls nicht.

\subsection{Ein generatives Marktmodell statt Rolling Windows}

Die zweite Idee stammt aus der Paper-Auswahl (Kuo et al., LOB-GAN):
Generalisiert der Agent besser, wenn er \emph{mehrere plausible Marktpfade} sieht statt
eines einzigen historischen? Der naheliegende Weg wären rollierende Trainingsfenster ---
und der ist in diesem Protokoll \textbf{nicht baubar}. Die Umgebung erzwingt pro
Environment mehr als $2 \cdot 252 = 504$ Beobachtungen (\texttt{assert len(df) >
2*trading\_days}), und \texttt{trading\_days} steuert gleichzeitig die Episodenlänge von
252 Schritten; das Minimum lässt sich also nicht absenken, ohne das Protokoll zu brechen.
Das Trainingsfenster hat aber nur \textbf{634} nutzbare Zeilen, und die Kursreihe beginnt
exakt am Fensteranfang. Zwei echte Sub-Fenster bräuchten je rund 550 Zeilen --- unmöglich.
Der Code prüft das inzwischen selbst und bricht mit klarer Meldung ab.

Stattdessen erzeugt ein \textbf{Block-Bootstrap} synthetische Kursreihen. Entscheidend ist
dabei, \emph{was} resampelt wird: Preisniveaus wären falsch, weil an jeder Blocknaht ein
künstlicher Kurssprung entstünde (Block endet bei 100, der nächste beginnt bei 80 ---
minus 20\,\% über Nacht). Gebootstrappt wird deshalb die \emph{Tagesform}:
$\log(\text{Close}_t/\text{Close}_{t-1})$, $\log(\text{Open}_t/\text{Close}_{t-1})$,
$\log(\text{High}_t/\text{Close}_t)$, $\log(\text{Low}_t/\text{Close}_t)$ und
$\log(\text{Volume}_t)$. Blöcke von 20 aufeinanderfolgenden Tagen werden mit Zurücklegen
gezogen und der Pfad wird vom ersten echten Close aus neu aufgebaut. Innerhalb der Blöcke
bleibt die Autokorrelation erhalten, die Kurve ist per Konstruktion stetig, und der Agent
sieht den \emph{realen} Trainingspfad nie. Gegenprobe: Die Streuung der Tagesrenditen
beträgt in den synthetischen Reihen $0{,}0229$ bzw.\ $0{,}0247$ gegenüber $0{,}0239$ in
der echten Reihe.

\subsection{Isolation: der Bootstrap trägt, die zwei Passagiere nicht}

In der Zeile ``Bündel $+$ Bootstrap'' steckt der Bootstrap noch zusammen mit den zwei
inerten Faktoren. Der isolierende Lauf \texttt{nstep} $10$ $+$ Bootstrap ohne
\texttt{ent\_coef} und \texttt{obs\_noise} schließt die Lücke (letzte Zeile):

\begin{table}[H]
  \centering
  \caption{Gepaarte Tests über die 20 Seed-Differenzen (df$=19$), Testmarge.}
  \label{tab:a3c-paired}
  \small
  \begin{tabular}{lrrr}
    \toprule
    Vergleich & $\Delta$ Testmarge & $t$ & $p$\\
    \midrule
    \texttt{nstep} 10 $-$ Baseline            & $+0{,}107$ & $1{,}69$ & $0{,}107$\\
    \texttt{ent\_coef} 0,02 $-$ Baseline      & $+0{,}028$ & $0{,}47$ & $0{,}643$\\
    \texttt{obs\_noise} 0,01 $-$ Baseline     & $+0{,}020$ & $0{,}36$ & $0{,}720$\\
    Bündel $-$ Baseline                       & $+0{,}126$ & $1{,}98$ & $0{,}062$\\
    Bündel $-$ \texttt{nstep} 10              & $+0{,}020$ & $0{,}24$ & $0{,}810$\\
    \texttt{nstep} 10 $+$ Bootstrap $-$ Baseline & $+0{,}200$ & $2{,}35$ & $\mathbf{0{,}030}$\\
    \texttt{nstep} 10 $+$ Bootstrap $-$ \texttt{nstep} 10 & $+0{,}093$ & $1{,}16$ & $0{,}260$\\
    \texttt{nstep} 10 $+$ Bootstrap $-$ Bündel $+$ Bootstrap & $+0{,}0002$ & $0{,}003$ & $0{,}998$\\
    \bottomrule
  \end{tabular}
\end{table}

\paragraph{Der Bootstrap-Effekt hält isoliert --- und war nie von den Passagieren getragen.}
Der isolierte Zuwachs beträgt $+0{,}093$ und damit \emph{exakt} denselben Wert wie die
kontaminierte Schätzung ($+0{,}093$); das Verhältnis ist $1{,}00$. Der Effekt ist also
weder geschrumpft noch verschwunden. Einschränkung: Für sich allein ist er
\textbf{nicht signifikant} ($p=0{,}26$, 12 von 20 Seeds positiv, Median $+0{,}138$,
getrimmter Mittelwert $+0{,}112$) --- der Punktschätzer ist stabil, die Stichprobe zu klein.

\paragraph{\texttt{nstep} 10 $+$ Bootstrap ist von ``Bündel $+$ Bootstrap'' nicht zu
unterscheiden} ($+0{,}0002$, $p=0{,}998$). Damit sind \texttt{ent\_coef} und
\texttt{obs\_noise} \textbf{endgültig als Passagiere belegt}: Die beiden Läufe
unterscheiden sich nur in diesen zwei Parametern und liefern praktisch dieselbe
Verteilung. (Nebenbei ein Konsistenzbeleg dafür, dass die Bootstrap-Replikate
deterministisch sind.)

\paragraph{Erstmals ein signifikanter Vorsprung.} \texttt{nstep} $10$ $+$ Bootstrap ist
gegenüber der Baseline signifikant besser ($+0{,}200$, $p=0{,}030$) und die eigene
Testmarge ist klar von null verschieden ($p=0{,}0010$, zweiseitig $0{,}0019$). Das ist der
\emph{erste} Arm in diesem Projekt, der die 5-\%-Schwelle im gepaarten Test gegen die
Baseline reißt --- weder \texttt{nstep} allein ($p=0{,}107$) noch das Bündel
($p=0{,}062$) schaffen das.

\begin{figure}[H]
  \centering
  \includegraphics[width=\textwidth]{update_2026_a3c.png}
  \caption{Ein-Faktor-Zerlegung und Block-Bootstrap, 20 Seeds. Links die Testmarge
  (Fehlerbalken: 95-\%-KI), rechts die Marge im Trainingsfenster. Die beiden
  Bootstrap-Arme (dunkel) haben in-sample praktisch null Vorsprung und out-of-sample den
  größten.}
  \label{fig:a3c}
\end{figure}

\subsection{Was das Overfitting-Bild zeigt}

Abbildung~\ref{fig:a3c} macht den Kern sichtbar. Die vier Arme, die auf dem echten
Fenster trainieren, tragen in-sample zwischen $+6{,}7$ und $+8{,}8$ NAV-Punkte --- und
davon bleibt out-of-sample fast nichts. Die beiden Bootstrap-Arme haben in-sample
$+0{,}02$, also \emph{null} Vorsprung, und liefern gleichzeitig die besten Testwerte. Sie
haben den historischen Pfad nie gesehen; dass sie im Trainingsfenster nicht gewinnen, ist
deshalb keine Schwäche, sondern der Mechanismus. Ein In-Sample-Vorsprung von $+8$ ist
damit als das entlarvt, was er ist: Memorieren.

\subsection{Einschränkungen}

\paragraph{20 Seeds sind eine kleine Stichprobe.} Die Konfidenzintervalle in
Tabelle~\ref{tab:a3c-arms} sind breit, und der zentrale Bootstrap-Zuwachs ist mit
$p=0{,}26$ unterpowert. Belastbar ist: dass die Baseline keinen Vorsprung hat, dass
\texttt{nstep} $10$ das Bündel erklärt, dass die zwei übrigen Faktoren nichts tun, und
dass \texttt{nstep} $10$ $+$ Bootstrap die Baseline signifikant schlägt. Nicht belastbar
ist die Behauptung, der Bootstrap-Zusatz sei \emph{für sich} signifikant.

\paragraph{Nur fünf verschiedene synthetische Pfade.} Bei 20 Agenten und fünf
Bootstrap-Replikaten teilen sich je vier Seeds dieselbe synthetische Reihe. Die Seeds
variieren dann nur die Exploration, nicht die Umgebung. Ein Lauf mit einem eigenen Pfad
pro Agent wäre der nächste saubere Schritt.

\paragraph{Die In-Sample-Marge der Bootstrap-Arme ist nicht vergleichbar.} Sie wird auf
dem echten Fenster gemessen, auf dem diese Agenten nie trainiert haben; die Zeile ist
deshalb nur als Overfitting-Indikator lesbar, nicht als Leistungsmaß.

\paragraph{Daten und Umgebung.} Es gelten die Einschränkungen aus
Abschnitt~\ref{sec:abweichungen}: Beobachtungsvektor 334 statt 205, lokale CSVs statt
\texttt{yfinance}. A3C nutzt das gleiche Environment wie A2C, die Abweichungen wirken also
auf beide Arme gleichermaßen.

\subsection{Fazit des A3C-Teils}

Die Baseline bestätigt den A2C-Befund auf einem zweiten Algorithmus: kein belegbarer
Vorsprung vor Buy-and-Hold. Die Zerlegung zeigt, dass der Effekt der naheliegenden
``Regularisierung'' in Wahrheit aus der \textbf{Kredit-Zuweisung} stammt
(\texttt{nstep} $5\to10$); Entropie-Bonus und Beobachtungsrauschen sind in dieser
Größenordnung wirkungslos. Der \textbf{Block-Bootstrap} liefert als zweiter Hebel den
besten Testwert und den ersten signifikanten Vorsprung gegenüber der Baseline ---
allerdings erst in der Kombination. Konzeptionell ist das der interessantere Befund: Der
Agent braucht den realen historischen Pfad nicht, um auf dem Testfenster gleichauf oder
besser zu sein; viele plausible Pfade genügen.
